\chapter{Lesson 5: Numbers, Dates and Time}

\section{Numbers}

\subsection{0 to 99}

Chinese numbers are completely regular: 11 is ``ten-one'', 20 is ``two-ten'', 35 is
``three-ten-five''.

\begin{center}
\begin{tabular}{llllll}
\toprule
0 & \zh{零} líng & 10 & \zh{十} shí & 20 & \zh{二十} èrshí \\
1 & \zh{一} yī & 11 & \zh{十一} shíyī & 21 & \zh{二十一} èrshíyī \\
2 & \zh{二} èr & 12 & \zh{十二} shí'èr & 35 & \zh{三十五} sānshíwǔ \\
3 & \zh{三} sān & 13 & \zh{十三} shísān & 48 & \zh{四十八} sìshíbā \\
4 & \zh{四} sì & 14 & \zh{十四} shísì & 60 & \zh{六十} liùshí \\
5 & \zh{五} wǔ & 15 & \zh{十五} shíwǔ & 77 & \zh{七十七} qīshíqī \\
6 & \zh{六} liù & 16 & \zh{十六} shíliù & 99 & \zh{九十九} jiǔshíjiǔ \\
7 & \zh{七} qī & 17 & \zh{十七} shíqī & & \\
8 & \zh{八} bā & 18 & \zh{十八} shíbā & & \\
9 & \zh{九} jiǔ & 19 & \zh{十九} shíjiǔ & & \\
\bottomrule
\end{tabular}
\end{center}

\subsection{Larger numbers}

\begin{itemize}
  \item \zh{百} bǎi = 100: \zh{一百二十} 120, \zh{三百零五} 305 (an internal zero is
        read as \zh{零}).
  \item \zh{千} qiān = 1000: \zh{两千零一十六} 2016.
  \item \zh{万} wàn = 10\,000. Chinese groups digits in fours, not threes:
        \zh{三万五千} = 35\,000; \zh{十万} = 100\,000.
\end{itemize}

\begin{aside}
Phone numbers and room numbers are read digit by digit, and \zh{一} is often said
\emph{yāo} to avoid confusion with \zh{七} qī: room 107 is \zh{幺零七} \emph{yāo líng
qī}. Hand signs for 6--10 differ from North American counting; 6 is the ``shaka'' sign
and 10 is crossed index fingers or a fist.
\end{aside}

\subsection{\zh{多少} vs \zh{几}}

\zh{多少} \emph{duōshao} ``how many/how much'' is used when the answer may be large; the
measure word is optional. \zh{几} expects a small number and needs a measure word.
\begin{itemize}
  \item \zh{你们班有多少(个)学生？} How many students are in your class?
  \item \zh{你的电话号码是多少？} What is your phone number?
  \item \zh{你有几个中国朋友？} How many Chinese friends do you have?
\end{itemize}

\section{Dates}

\subsection{Months, days and weekdays}

Months are simply numbered: \zh{一月} January, \zh{二月} February, \ldots,
\zh{十二月} December. Days of the month use \zh{号} hào in speech and \zh{日} rì in
writing. Weekdays are numbered from Monday.

\begin{center}
\begin{tabular}{llll}
\toprule
Characters & Pinyin & Meaning & Note \\
\midrule
\zh{星期一} & xīngqīyī & Monday & ``week-one'' \\
\zh{星期二} & xīngqī'èr & Tuesday & \\
\zh{星期三} & xīngqīsān & Wednesday & \\
\zh{星期四} & xīngqīsì & Thursday & \\
\zh{星期五} & xīngqīwǔ & Friday & \\
\zh{星期六} & xīngqīliù & Saturday & \\
\zh{星期天}/\zh{星期日} & xīngqītiān / xīngqīrì & Sunday & \zh{天} in speech \\
\zh{今天} & jīntiān & today & \\
\zh{明天} & míngtiān & tomorrow & \\
\zh{昨天} & zuótiān & yesterday & \\
\zh{今年} & jīnnián & this year & \\
\zh{生日} & shēngrì & birthday & \\
\bottomrule
\end{tabular}
\end{center}

\subsection{Order: from largest to smallest}

\begin{important}[Big before small]
Chinese orders dates (and times, and addresses) from the largest unit to the smallest:
year, month, day, weekday, part of day, hour, minute.
\end{important}

\begin{center}
\begin{tikzpicture}[u/.style={draw, rounded corners, minimum height=0.9cm, align=center, font=\small}]
  \node[u, fill=blue!25] (y) at (0,0) {\zh{二〇一六年}\\ 2016};
  \node[u, fill=blue!20] (m) at (2.4,0) {\zh{九月}\\ Sept.};
  \node[u, fill=blue!15] (d) at (4.3,0) {\zh{八号}\\ 8th};
  \node[u, fill=blue!10] (w) at (6.3,0) {\zh{星期四}\\ Thu.};
  \node[u, fill=blue!6] (p) at (8.3,0) {\zh{上午}\\ a.m.};
  \node[u, fill=blue!3] (h) at (10.4,0) {\zh{十点半}\\ 10:30};
  \foreach \a/\b in {y/m,m/d,d/w,w/p,p/h} \draw[->, thick] (\a) -- (\b);
  \draw[|-{Stealth}, gray] (0,-0.95) -- (10.4,-0.95) node[midway, below, font=\scriptsize] {larger unit $\rightarrow$ smaller unit};
\end{tikzpicture}
\end{center}

Years are read digit by digit: 2016 = \zh{二〇一六年} \emph{èr líng yī liù nián}
(\zh{〇} is the written form of zero in dates).

\subsection{Nominal predicates}

When telling dates, times and ages, \zh{是} is usually left out: a noun phrase acts as
the predicate. In the negative, \zh{是} must come back.
\begin{itemize}
  \item \zh{今天九月十二号。} Today is September 12.
  \item \zh{明天星期五。} Tomorrow is Friday.
  \item \zh{今天不是星期五，是星期四。} Today isn't Friday, it's Thursday.
  \item \zh{我今年十九岁。} I'm 19 this year.
\end{itemize}

\subsection{Asking the date and age}

\begin{itemize}
  \item \zh{今天几月几号？} What's the date today?
  \item \zh{今天星期几？} What day of the week is it?
  \item \zh{你的生日是几月几号？} When is your birthday?
  \item \zh{你今年多大？} How old are you this year? (to peers)
  \item \zh{你几岁？} How old are you? (to children)
  \item \zh{您多大年纪？} How old are you? (to the elderly; \zh{年纪} niánjì ``age'')
\end{itemize}

\begin{aside}
Chinese people often state their age by zodiac animal: \zh{我属狗。} \emph{Wǒ shǔ gǒu.}
``I was born in the Year of the Dog.'' Since the animals cycle every 12 years, this
tells the listener your age without asking directly. (The full zodiac comes back in
CHINA 102R.)
\end{aside}

\section{Telling time}

\begin{center}
\begin{tabular}{llll}
\toprule
Characters & Pinyin & Meaning & Example \\
\midrule
\zh{点(钟)} & diǎn(zhōng) & o'clock & \zh{三点} 3:00 \\
\zh{分} & fēn & minute & \zh{三点十分} 3:10 \\
\zh{刻} & kè & quarter hour & \zh{三点一刻} 3:15 \\
\zh{半} & bàn & half & \zh{三点半} 3:30 \\
\zh{差} & chà & lacking, to & \zh{差五分四点} 3:55 \\
\zh{现在} & xiànzài & now & \zh{现在几点？} \\
\zh{早上} & zǎoshang & morning & \\
\zh{上午} & shàngwǔ & morning (before noon) & \\
\zh{中午} & zhōngwǔ & noon & \\
\zh{下午} & xiàwǔ & afternoon & \\
\zh{晚上} & wǎnshang & evening & \\
\bottomrule
\end{tabular}
\end{center}

\begin{center}
\begin{tikzpicture}[scale=1.15]
  \draw[thick] (0,0) circle (1.6);
  \foreach \h in {1,...,12} {
    \draw (90-\h*30:1.45) -- (90-\h*30:1.6);
    \node[font=\scriptsize] at (90-\h*30:1.25) {\h};
  }
  % 3:15 -> minute hand to 3 (0 deg), hour hand a bit past 3
  \draw[very thick, -{Stealth}] (0,0) -- (0:1.25);
  \draw[line width=2pt, -{Stealth}] (0,0) -- (-7.5:0.8);
  \fill (0,0) circle (0.05);
  \node[below, font=\small] at (0,-1.7) {\zh{三点一刻} 3:15};
  \begin{scope}[xshift=4.5cm]
  \draw[thick] (0,0) circle (1.6);
  \foreach \h in {1,...,12} {
    \draw (90-\h*30:1.45) -- (90-\h*30:1.6);
    \node[font=\scriptsize] at (90-\h*30:1.25) {\h};
  }
  \draw[very thick, -{Stealth}] (0,0) -- (-90:1.25);
  \draw[line width=2pt, -{Stealth}] (0,0) -- (90-6.5*30:0.8);
  \fill (0,0) circle (0.05);
  \node[below, font=\small] at (0,-1.7) {\zh{六点半} 6:30};
  \end{scope}
  \begin{scope}[xshift=9cm]
  \draw[thick] (0,0) circle (1.6);
  \foreach \h in {1,...,12} {
    \draw (90-\h*30:1.45) -- (90-\h*30:1.6);
    \node[font=\scriptsize] at (90-\h*30:1.25) {\h};
  }
  \draw[very thick, -{Stealth}] (0,0) -- (90-55*6:1.25);
  \draw[line width=2pt, -{Stealth}] (0,0) -- (90-9.9*30:0.8);
  \fill (0,0) circle (0.05);
  \node[below, font=\small] at (0,-1.7) {\zh{差五分十点} 9:55};
  \end{scope}
\end{tikzpicture}
\end{center}

Time expressions also go big-to-small: \zh{明天晚上八点} ``8 p.m.\ tomorrow''.

\section{Time--place--action word order}

\begin{important}[Time before the verb]
A time-when expression comes \emph{before} the verb, either right after the subject or
at the very start of the sentence. It never goes at the end as in English.
\begin{itemize}
  \item \zh{我明天晚上请你吃饭。} I'll treat you to dinner tomorrow night.
  \item \zh{明天晚上我请你吃饭。} (same meaning; topic first)
  \item *\zh{我请你吃饭明天晚上。} (wrong)
\end{itemize}
\end{important}

\section{Making a birthday invitation}

\begin{center}
\begin{tabular}{llll}
\toprule
Characters & Pinyin & Part of speech & Meaning \\
\midrule
\zh{请客} & qǐngkè & verb-object & to treat (pay for others) \\
\zh{吃饭} & chīfàn & verb-object & to eat a meal \\
\zh{怎么样} & zěnmeyàng & question word & how about it \\
\zh{太\ldots{}了} & tài \ldots le & & too / extremely \ldots \\
\zh{为什么} & wèishénme & question word & why \\
\zh{因为} & yīnwèi & conjunction & because \\
\zh{还是} & háishi & conjunction & or (in questions) \\
\zh{可是} & kěshì & conjunction & but \\
\zh{喜欢} & xǐhuan & verb & to like \\
\zh{菜} & cài & noun & dish, cuisine \\
\zh{认识} & rènshi & verb & to know (a person) \\
\zh{朋友} & péngyou & noun & friend \\
\bottomrule
\end{tabular}
\end{center}

\subsection{Suggestions with \zh{怎么样}}

``Proposal, \zh{怎么样}？'' = ``How about \ldots?''
\begin{itemize}
  \item \zh{星期六我请你吃饭，怎么样？}---\zh{太好了！谢谢！}
\end{itemize}

\subsection{Alternative questions with \zh{还是}}

``A \zh{还是} B'' asks the listener to choose; no \zh{吗}.
\begin{itemize}
  \item \zh{你喜欢吃中国菜还是美国菜？} Do you like Chinese food or American food?
  \item \zh{你是大学生还是研究生？} Are you an undergraduate or a graduate student?
\end{itemize}

\begin{example}[Short exchange]\leavevmode
\begin{itemize}
  \item \textbf{A:} \zh{小林，九月十二号是星期几？}
  \item \textbf{B:} \zh{是星期四。为什么？}
  \item \textbf{A:} \zh{那天是我的生日。我请你吃饭，怎么样？}
  \item \textbf{B:} \zh{太好了！几点？}
  \item \textbf{A:} \zh{晚上七点半。}
\end{itemize}
\end{example}

\section*{Review}
\begin{itemize}
  \item Numbers are regular; \zh{万} = 10\,000 groups digits in fours.
  \item \zh{多少} for possibly large numbers, \zh{几} + MW for small ones.
  \item Dates and times go from big to small: year, month, day, weekday, part of day,
        hour, minute.
  \item Dates, times and ages can be predicates without \zh{是}; the negative needs
        \zh{不是}.
  \item Time: \zh{点, 分, 刻, 半, 差}.
  \item Time-when expressions go before the verb.
  \item \zh{怎么样} for suggestions; A \zh{还是} B for choices.
\end{itemize}
